% Fuente: Pauta del Auxiliar 3, «Convexidad — Soluciones de una inecuación cuadrática».
{\color{MidnightBlue}
\begin{itemize}
  \item[(a)]
  \begin{itemize}
  \item[$\Rightarrow$:] Supongamos que $S$ es convexo. Queremos probar que $S\cap L$ es convexo.

  Sea arbitraria la línea 
  $L=\{\hat x+t v\mid t\in\mathbb{R}\}$ y tomemos dos puntos 
  \[
    y_1 = \hat x + t_1 v,\quad 
    y_2 = \hat x + t_2 v
  \]
  en $S\cap L$. Para cualquier $\lambda\in[0,1]$ tenemos
  \[
    \lambda y_1 + (1-\lambda)y_2 
    = \lambda(\hat x + t_1 v) + (1-\lambda)(\hat x + t_2 v)
    = \hat x + (\lambda t_1 + (1-\lambda)t_2)\,v,
  \]
  que pertenece a $L$. Como además $S$ es convexo, $\lambda y_1+(1-\lambda)y_2\in S$.  
  Por lo tanto $\lambda y_1 + (1-\lambda)y_2 \in S\cap L$, y $S\cap L$ es convexo.
  \item[$\Leftarrow$:] Ahora supongamos que para toda línea $L$ la intersección $S\cap L$ es convexa. 
  Queremos probar que $S$ es convexo en $\mathbb{R}^n$. 

  Tomemos dos puntos arbitrarios $x_1,x_2\in S$. Definimos la línea que los contiene:
  \[
    L = \{\,x_1 + t\,(x_2 - x_1)\mid t\in\mathbb{R}\,\}.
  \]
  Como $x_1,x_2\in S\cap L$ y por hipótesis $S\cap L$ es convexo, para todo $\lambda\in[0,1]$ se cumple
  \[
    \lambda x_1 + (1-\lambda)x_2 
    \;\in\; S\cap L
    \;\subseteq\; S.
  \]
  Esto demuestra que $S$ es convexo.
  \end{itemize}
  \item[(b)] Suponga que
    \[
    A\succeq 0.
    \]
    Queremos demostrar que \(C\) es convexo.

    Por el resultado anterior, basta estudiar la intersección de \(C\) con una línea arbitraria. Sea
    \[
    x=\hat{x}+tv,
    \]
    donde \(\hat{x},v\in\mathbb{R}^n\) y \(t\in\mathbb{R}\).

    Reemplazando en la desigualdad cuadrática:
    \[
    (\hat{x}+tv)^\top A(\hat{x}+tv)
    +
    b^\top(\hat{x}+tv)
    +
    c
    \leq 0.
    \]

    Expandiendo:
    \[
    \hat{x}^\top A\hat{x}
    +
    2t v^\top A\hat{x}
    +
    t^2 v^\top A v
    +
    b^\top \hat{x}
    +
    t b^\top v
    +
    c
    \leq 0.
    \]

    Agrupando en términos de \(t\):
    \[
    \phi(t)
    =
    (v^\top A v)t^2
    +
    \left(2v^\top A\hat{x}+b^\top v\right)t
    +
    \hat{x}^\top A\hat{x}
    +
    b^\top \hat{x}
    +
    c
    \leq 0.
    \]

    Como \(A\succeq 0\), se cumple que
    \[
    v^\top A v\geq 0
    \qquad
    \forall v\in\mathbb{R}^n.
    \]
    Por lo tanto, \(\phi(t)\) es una función convexa en \(t\).

    Luego, el conjunto
    \[
    \{t\in\mathbb{R}\mid \phi(t)\leq 0\}
    \]
    es un subnivel \footnote{es un subconjunto de $\mathbb{R}$ formado por todos los valores de t que cumplen la condición $ϕ(t)≤0$} de una función convexa de una variable, y por tanto es convexo.

    Como esto ocurre para toda línea, se concluye que \(C\) es convexo.

    \[
    A\succeq 0 \Rightarrow C\text{ es convexo.}
    \]
\end{itemize}
}
